\documentclass[../Main.tex]{subfiles}

\begin{document}
\section{Fourier Transforms}
\subsection{Defining the Fourier Transform}
\begin{definition}{Absolute convergence}
    A function $f : \R\mapsto\C$ is \underline{absolutely convergent} if:
    \begin{equation*}
        \int_{-\infty}^{\infty} \abs{f(x)} dx \text{ is finite.}
    \end{equation*}
\end{definition}
\begin{definition}{Fourier transform}
    Let $f : \R \mapsto \C$ be absolutely convergent, have a bounded variation, and a finite number of discontinuities. Then its \underline{Fourier transform} (and its inverse) are defined as:
    \begin{equation*}
        \ft{f}(k) = \FT[f(x)](k) = \int_{-\infty}^{\infty} f(x) e^{-ikx} dx
    \end{equation*}
    and:
    \begin{equation*}
        f(x) = \FT^{-1}[\ft{f}(k)](x) = \frac{1}{2\pi} \int_{-\infty}^{\infty} \ft{f}(k)e^{ikx} dk 
    \end{equation*}
\end{definition}
\begin{remarks}
    \item When possible, $\ft{f}(k)$ will be the Fourier transform of $f(x)$.
    \item There are many conventions for signs of $k$ and factors of $2\pi$.
    \item In physics, the argument to the Fourier transform has important meaning. If $x$ is a spatial variable, $k$ is a wave number. If $t$ is a temporal variable, $k$ is instead written $\omega$ and is the angular frequency.
    \item Some non-square integrable functions can be handled via distributions:
        \begin{equation*}
            f(x) = 1 \implies \ft{f}(k) = \int_{-\infty}^{\infty} e^{ikx} dx = 2\pi \delta(k)
        \end{equation*}
    \item At discontinuities, the Fourier transform returns the average,
        \begin{equation*}
            \frac{f(x^+) + f(x^-)}{2}
        \end{equation*}
\end{remarks}
\begin{definition}{Cauchy Principal Value}
    The \underline{Cauchy principal value} of an integral
    \begin{equation*}
        \int_{-\infty}^{\infty} g(x) dx
    \end{equation*}
    is:
    \begin{equation*}
        \dashint_{-\infty}^\infty g(x) dx = \lim_{R \to \infty} \int_{-R}^R g(x) dx
    \end{equation*}
\end{definition}
This allows us to give a value to integrals that would otherwise not be well-defined. For example,
\begin{equation*}
    \int_{-\infty}^{\infty} \frac{x}{1+x^2} dx \text{ is not well-defined}
\end{equation*}
but,
\begin{equation*}
    \dashint_{-\infty}^{\infty} \frac{x}{1+x^2} dx = 0
\end{equation*}
because we can cancel the positive and negative parts.

If we have a pole in the middle of the integration domain (here $\R$), we can also assign a Cauchy Principal Value in certain cases. In general,
\begin{equation*}
    \frac12 \left[\ft{f}(k^+) + \ft{f}(k^-)\right] = \dashint_{-\infty}^{\infty} f(x) e^{-ikx} dx 
\end{equation*}
\subsection{Fourier Transforms of Example Functions}
\begin{example}[Gaussian function]
    Consider:
    \begin{equation*}
        f(x) = e^{-x^2 / 2}
    \end{equation*}
    Then we find the Fourier Transform:
    \begin{align*}
        \ft{f}(k) &= \int_{-\infty}^{\infty} e^{-x^2 / 2 - ikx} dx \\
        &= \int_{-\infty}^{\infty} \exp\left[-\frac{(x+ik)^2}{2} - \frac{k^2}{2}\right] dx \\
        \intertext{Now let $z = x + ik$:}
        &= e^{-k^2 / 2} \int_{-\infty + ik}^{\infty + ik} e^{-z^2 / 2} dz
    \end{align*}
    Let this be the integral over the contour $\gamma_0$, which is a horizontal line given by $\Im(z) = ik$. We introduce $\gamma_R$ and $\gamma_{-R}$ to close the contour to $\gamma_1$ which is $\R$ traversed from $\infty$ to $-\infty$. Therefore the contour $\gamma = \gamma_0 + \gamma_R + \gamma_1 + \gamma_{-R}$ is a rectangular contour. We see that the integrals over $\gamma_{\pm R}$ must be zero as $R \to \infty$, and since there are no singularities inside the contour, \thmref{thmCauchy} tells us that:
    \begin{equation*}
        \int_{\gamma_0} e^{-z^2 / 2} dz = -\int_{\gamma_1} e^{-z^2 / 2} dz = \int_{-\infty}^{\infty} e^{-x^2} dx = \sqrt{2\pi}
    \end{equation*}
    and so:
    \begin{equation*}
        \ft{f}(k) = \sqrt{2\pi} e^{-k^2 / 2}
    \end{equation*}
    which is interesting, and tells us that Gaussians are eigenvectors of $\FT$.
\end{example}
\begin{example}[Exponentials]
    Consider:
    \begin{equation*}
        \ft{f}(k) = \frac{1}{a + ik},\quad a>0
    \end{equation*}
    then we note that $\ft{f} \to 0$ as $|k| \to \infty$, and so we can use \thmref{thmJordanLemma}.

    Consider:
    \begin{equation*}
        f(x) = \frac{1}{2\pi} \int_{-\infty}^{\infty} \ft{f}(k) e^{ikx} dk
    \end{equation*}
    if $x > 0$, we close the contour in the upper half-plane using $\gamma_R$ and find:
    \begin{equation*}
        \lim_{R \to \infty} \int_{\gamma_R} \ft{f}(k) dk = 0
    \end{equation*}
    now we need the residue at the pole:
    \begin{equation*}
        \res_{k = ia} \frac{e^{ikx}}{a + ik} -ie^{ax}
    \end{equation*}
    and so, for $x > 0$,
    \begin{equation*}
        f(x) = e^{-ax}
    \end{equation*}
    But since there are no poles in the lower half-plane, when we close the contour there we get that the integral over it is zero. \Thmref{thmJordanLemma} tells us that the integral over the semicircle is zero, so also the integral over the real axis must be zero. Therefore,

    \begin{equation*}
        f(x) = 
        \begin{cases}
            e^{-ax} & x > 0 \\
            0 & x < 0
        \end{cases}
    \end{equation*}
    We can check our answer,
    \begin{align*}
        \ft{f}(k) &= \int_{-\infty}^{\infty} f(x)e^{-ikx} dx \\
        &= \int_{0}^{\infty} e^{-ax - ikx} dx \\
        &= \frac{1}{a + ik} \text{ as required.}
    \end{align*}
\end{example}
\section{Laplace Transforms}
This section is motivated by the desire to compute something akin to a Fourier transform for functions that diverge as $t \to \infty$.
\subsection{Defining the Laplace Transform}
\begin{definition}{Laplace Transform}
    Let $f(t)$ be a function defined for $t > 0$. The \underline{Laplace transform} of $f(t)$ is given by:
    \begin{equation}
        F(s) = \LT[f](s) = \int_{0}^{\infty} f(t)e^{-st} dt
        \label{eqnLaplaceTransform}
    \end{equation}
    provided that the integral exists.
\end{definition}
\begin{remarks}
    \item This is valid as long as the function $f(t)$ grows slower than exponential. For example, any polynomial.
    \item Sometimes $p$ is used instead of $s$ as the argument to the Laplace Transform. We will use $s$, but this is good to know.
    \item If $f(t) = 0$ for $t < 0$, then we have the following relation to the Fourier Transform:
        \begin{equation*}
            F(s) = \ft{f}(-is)
        \end{equation*}
\end{remarks}
\subsection{Laplace Transforms of Example Functions}
\begin{example}[Constant function]
    Let $f(t) = 1$.
    \begin{equation*}
        \int_{0}^{\infty} e^{-st} dt = \frac{1}{s}
    \end{equation*}
    What is going on? If we have $\Re(s) < 0$ then the integral should diverge. However, the Laplace Transform displays no such divergence for $\Re(s) < 0$, only at $s = 0$. We solve this by analytic continuation, where we define the Laplace Transform for $\Re(s) > 0$ to be $\frac1s$, and we define it to be the same thing on $\Re(s) < 0$.
\end{example}
\begin{example}[Linear function]
    Consider $f(t) = t$.
    \begin{align*}
        F(s) = \int_{0}^{\infty} te^{-st} dt \\
        &= \left[t \left(-\frac1s\right)e^{-st}\right]_0^\infty - \int_{0}^{\infty} \left(-\frac1s\right)e^{-st} dt \\
        &= \frac{1}{s^2}
    \end{align*}
    and again we argue by analytic continuation that we can define this on $\C \backslash \{0\}$.
\end{example}
\begin{example}[Exponentials]
    For constant $\lambda$, we can formally consider:
    \begin{equation*}
        \LT\left[e^{\lambda t}\right](s) = \int_{0}^{\infty} e^{-(s - \lambda)t} dt = \frac{1}{s - \lambda}
    \end{equation*}
    and by analytic continuation, we can extend this beyond $\Re(s - \lambda) > 0$ to $\C \backslash \{\lambda\}$.
\end{example}
\begin{example}[Sine]
    Consider:
    \begin{align*}
        \LT[\sin](s) &= \LT\left[\frac1{2i} (e^{it} - e^{-it})\right](s) \\
        &= \frac{1}{2i} \left(\frac{1}{s-i} - \frac{1}{s+i}\right) = \frac{1}{s^2 + 1}
    \end{align*}
    we have just required that the Laplace transform be a linear map. In the following section, we will prove properties about the Laplace transform including linearity.
\end{example}
\subsection{Properties of the Laplace Transform}
\begin{propositions}{
        Let $f, g : \C \mapsto \C$ with Laplace transforms $F$ and $G$, and $\alpha, \beta \in \C$.
        \label{propsLTProps}
    }
    \item Linearity:
        \begin{equation*}
            \LT\left[\alpha f + \beta g\right] = \alpha \LT[f] + \beta \LT[g]
        \end{equation*}
    \item Translation: let $t_0 \in \R$, then
        \begin{equation*}
            \LT[f(t - t_0) H(t - t_0)] (s) = e^{-s t_0} F(s)
        \end{equation*}
        where $H$ is the Heaviside function.
    \item Scaling: Let $\lambda \in \R$ be positive, then
        \begin{equation*}
            \LT[f(\lambda t)] = \frac{1}{\lambda} F(s / \lambda)
        \end{equation*}
    \item Shifting by a constant:
        \begin{equation*}
            \LT[e^{s_0 t} f(t)] = F(s - s_0)
        \end{equation*}
    \item Derivatives:
        \begin{equation*}
            \LT[f'(t)] = sF(s) - f(0)
        \end{equation*}
        and the property that $f(0)$ appears in this term is known as \textit{memory}. The second derivative follows immediately:
        \begin{equation*}
            \LT[f''(t)] = s^2 F(s) - sf(0) - f'(0)
        \end{equation*}
    \item Derivative of the Laplace transform:
        \begin{equation*}
            \LT[tf(t)] = -F'(s)
        \end{equation*}
        this is useful because, given the Laplace Transform of $f(t)$, we can immediately know the Laplace transform of $tf(t)$.
    \item If $\lim_{t \to\infty} f(t)$ exists, then:
        \begin{equation*}
            \lim_{s \to \infty} sF(s) = f(0) \text{ and } \lim_{s \to 0} sF(s) = \lim_{t \to \infty}f(t)
        \end{equation*}
\end{propositions}
\begin{proof}
    \begin{enumerate}
        \item Simple from the definitions
        \item Set $\tilde{t} = t - t_0$, and consider:
            \begin{align*}
                \int_{0}^{\infty}& g(t - t_0) e^{-st} dt = \int_{-t_0}^{\infty} g(\tilde{t}) e^{-s(\tilde{t} + t_0)} d\tilde{t} \\
                &= e^{-st_0} \int_{-t_0}^{\infty} g(\tilde{t}) e^{-s\tilde{t}} d\tilde{t}
            \end{align*}
            and so now take $g(t) = f(t) H(t)$:
            \begin{align*}
                \int_{0}^{\infty}& f(t-t_0) H(t - t_0) e^{-st} dt = e^{-st_0} \int_{-t_0}^{\infty} f(\tilde{t}) H(\tilde{t}) e^{-s\tilde{t}} d\tilde{t} \\
                &= e^{-st_0} \int_{0}^{\infty} f(\tilde{t}) e^{-s\tilde{t}} d\tilde{t} \\
                &= e^{-st_0} F(s)
            \end{align*}
        \item We calculate it directly:
            \begin{align*}
                \LT[f(\lambda t)] &= \int_{0}^{\infty} f(\lambda t) e^{-st} dt \\
                &= \frac1\lambda \int_{0}^{\infty} f(\tilde{t}) e^{-\frac{s}{\lambda}\tilde{t}} d\tilde{t} \text{ by taking } \tilde{t} = \lambda t \\
                &= \frac1\lambda F\left(\frac{s}{\lambda}\right)
            \end{align*}
        \item TODO %TODO
        \item Again we calculate this directly:
            \begin{align*}
                \LT[f'(t)](s) &= \int_{0}^{\infty} f'(t)e^{-st} dt \\
                &= \left[f(t) e^{-st}\right]_0^\infty - \int_{0}^{\infty} (-s) f(t) e^{-st} dt \\
                &= sF(s) - f(0)
            \end{align*}
            where we are assuming that $f(t)$ does not grow faster than exponentials as $t \to \infty$.
        \item We again integrate by parts:
            \begin{align*}
                F'(s) &= \frac{d}{ds} \int_{0}^{\infty} f(t) e^{-st} dt \\
                &= \int_{0}^{\infty} (-t)f(t)e^{-st} dt \\
                &= \LT[-t f(t)]
            \end{align*}
        \item We have seen, in the derivative section, that:
            \begin{equation*}
                sF(s) = f(0) + \int_{0}^{\infty} f'(t) s^{st} dt
            \end{equation*}
            and so:
            \begin{align*}
                \lim_{s \to 0} sF(s) &= f(0) + \underbrace{\int_{0}^{\infty} f'(t) dt}_I \\
                &= \lim_{t \to \infty} f(t) - f(0) + f(0) = \lim_{t \to \infty} f(t).
            \end{align*}
            Also, we know that the limit exists, so $f(t)$ and $f'(t)$ do not grow faster than exponential. Then the integral $I$ vanishes as $s \to \infty$ and:
            \begin{equation*}
                \lim_{s \to \infty}sF(s) = f(0)
            \end{equation*}
    \end{enumerate}
\end{proof}
Now we have these properties, we can often calculate Laplace Transforms much more easily.
\begin{examples}{
        We can use the properties for the following examples:
    }
    \item $f(t) = t\sin(t)$. We already know the LT of $\sin$, and we know that $t\sin(t)$ is related to the derivative of the LT. Therefore:
        \begin{align*}
            \LT[t\sin(t)](s) &= -\frac{d}{ds} \LT[\sin(t)](s) \\
            &= -\frac{d}{ds} \left(\frac{1}{1 + s^2}\right) \\
            &= \frac{2s}{1 + s^2}
        \end{align*}
    \item We already know the Laplace transform of the constant function:
        \begin{equation*}
            \LT[1](s) = \frac{1}{s}
        \end{equation*}
        then in order to find the monomial function $f(t) = t^n$, we can take $n$ derivatives of the constant function's Laplace Transform:
        \begin{align*}
            \LT[t^n](s) &= (-1)^n \frac{d^{n}}{ds^{n}} \left(\frac1s\right) \\
            &= \frac{m!}{s^m}
        \end{align*}
    \item The Gamma function, that generalises $(n-1)!$:
        \begin{equation*}
            \Gamma(n) = \int_{0}^{\infty} e^{-t} t^{n-1} dt
        \end{equation*}
        is given by $\Gamma(n) = \LT[t^{n-1}](s = 1)$. Then we can use this to calculate noninteger values for $n$.
    \item Scaling the argument to a sin function: 
        \begin{align*}
            \LT[\sin(a t)] &= \frac1a \LT[\sin(t)]\left(\frac{s}{a}\right) \\
            &= \frac1a \frac{1}{(s / a)^2 + 1} \\
            &= \frac{a}{s^2 + a^2}
        \end{align*}
    \item We can find the cosine function using linearity, exponentials and sines:
        \begin{equation*}
            \LT[e^{iat}](s) = \frac{1}{s-ia} = \frac{s + ia}{s^2 + a^2}
        \end{equation*}
        and we know that this is a sum of $cos(at) + i\sin(at)$, so:
        \begin{equation*}
            \LT[\cos(at)] = \frac{s + ia}{s^2 + a^2} - \frac{ia}{s^2 + a^2} = \frac{s}{s^2 + a^2}
        \end{equation*}
        We could also have found this using the derivative.
\end{examples}
\subsection{Inverse Laplace Transform}
\begin{proposition}
    Given a function $F(s)$, we can compute the inverse Laplace transform $f(t)$ by:
    \begin{equation}
        f(t) = \frac1{2\pi i} \int_{\alpha - i \infty}^{\alpha + i \infty} F(s) e^{st} ds 
        \label{eqnInverseLaplace}
    \end{equation}
    and this is known as the \underline{Bromwich Inversion Formula}.

    The constant $\alpha$ is chosen such that the \underline{Bromwich inversion contour} $\gamma$ defined by $\{\Re(s) = \alpha\}$ lies to the right of all singularities of $F(s)$.
    \label{propInverseLaplace}
\end{proposition}
\begin{proof}
    Since $f(t)$ has a Laplace Transform, we have $f(t) = 0$ for $t < 0$, and $f$ does not grow faster than exponential.

    We can therefore choose an $\alpha \in \R$ such that:
    \begin{equation*}
        g(t) = f(t) e^{-\alpha t}
    \end{equation*}
    has a Fourier Transform (since it decays exponentially (or faster) as $t \to \infty$, and is zero for $t < 0$). This will be given by:
    \begin{equation*}
        \ft{g}(\omega) = \int_{-\infty}^{\infty} f(t)e^{\alpha t - i \omega t} dt
    \end{equation*}
    then we note that this looks like a Laplace Transform of $f$:
    \begin{equation*}
        \ft{g}(\omega) = \int_{-\infty}^{\infty} f(t)e^{-(\alpha + i \omega)t} dt = F(\alpha + i \omega)
    \end{equation*}
    which we can invert (by the Fourier Inversion Formula):
    \begin{equation*}
        g(t) = \frac{1}{2\pi} \int_{-\infty}^{\infty} F(\alpha + i \omega) e^{i \omega t} d\omega
    \end{equation*}
    and finally, take $s = \alpha + i \omega$:
    \begin{align*}
        f(t) e^{\alpha t} &= \frac{1}{2\pi i} \int_{\alpha - i \infty}^{\alpha + i \infty} F(s) e^{st - \alpha t} ds \\
        f(t) &= \int_{\alpha - i \infty}^{\alpha + i \infty} F(s) e^{st} ds
    \end{align*}
    The added condition that the contour lies to the right of all singularities of $F(s)$ fixes a constant of integration, and ensures that $f(t) = 0$ on $t < 0$. This is because we can close the contour $\gamma$ using a semicircle that avoids all singularities, meaning that the integral is zero.
\end{proof}
\end{document}